两能级系统哈密顿量 \(H=\hbar\Omega\sigma_x\)，\(\Omega>0\)，初态 \(|0\rangle\) 为 \(\sigma_z\) 正本征态。每隔 \(\tau=t/N\) 进行一次理想选择性投影测量 \(P=|0\rangle\langle0|\)，仅保留全部测量成功的轨迹。
(1) 求自由演化、生存概率和首次正交时间。
(2) 计算 \(N\) 次测量全部成功概率，严格求 \(N\to\infty\) 的固定时间极限与首阶误差。
(3) 对一般归一初态推导短时生存概率与能量方差关系，定义Zeno时间并讨论谱矩条件。
(4) 计算首次失败即停止时的期望测量次数，比较选择性与非选择性测量及非理想实验限制。
